% --- Archon physics formalization source begin ---
% archon:physics
% archon:covers PhyXMiniProblems/problem_phyx_mini_0612.lean
% archon:source-report reports/phyx_mini/problem_phyx_mini_0612.source.json

\chapter{Physics problem phyx\_mini\_0612}
\label{ch:PhyXMiniProblems_problem_phyx_mini_0612}

\paragraph{Problem source.}
\#\# Physical scenario

The work function of a material is examined using the apparatus shown schematically in figure. The ammeter measures the photocurrent. The voltage of the anode relative to the cathode is \$V\_\{AC\}\$.The switch is returned to position \$a\$, and the wavelength of the light is reduced to 65.0 nm. The photocurrent increases.

\#\# Question

What is the minimum value of \$R\$ for which the photocurrent will cease?

\#\# Answer choices

- A: A:  13.8kΩ

- B: B:  12.7kΩ

- C: C:  11.9kΩ

- D: D:  11.6kΩ

\#\# Figure caption (auxiliary; use the image as primary evidence)

The image shows an electric circuit diagram illustrating a photoelectric effect experiment setup. Here's a detailed description:

1. **Photoelectric Tube**:
   - On the left, there's a photoelectric tube with two parallel metal plates.
   - The left plate emits electrons (\textbackslash{}( e\textasciicircum{}- \textbackslash{})) when illuminated. For this, a photon (\textbackslash{}( \textbackslash{}lambda \textbackslash{})) hits the plate, causing an electron to be emitted with velocity (\textbackslash{}( v \textbackslash{})).
   - The symbol \textbackslash{}( \textbackslash{}phi \textbackslash{}) is present next to the emitted electron, likely representing the work function of the material.

2. **Voltage and Components**:
   - A voltage \textbackslash{}( -V\_\{AC\} \textbackslash{}) is applied between the plates (the left plate is grounded).
   - A resistor \textbackslash{}( R \textbackslash{}) is in series with the photoelectric tube.

3. **Current Measurement**:
   - An ammeter \textbackslash{}( A \textbackslash{}) is connected in series to measure the current \textbackslash{}( i \textbackslash{}), which is shown in the circuit with an arrow.

4. **Power Supply**:
   - Two batteries, each marked with \textbackslash{}( 15.0 \textbackslash{}, \textbackslash{}text\{V\} \textbackslash{}), are connected in series, providing a total of 30.0 V to the circuit.

5. **Resistor and Switch**:
   - A \textbackslash{}( 1.00 \textbackslash{}, \textbackslash{}text\{k\}\textbackslash{}Omega \textbackslash{}) resistor is present in the circuit.
   - A switch, labeled with \textbackslash{}( a \textbackslash{}) and \textbackslash{}( b \textbackslash{}), is included, allowing disconnection of part of the circuit.

6. **Current Flow**:
   - The current flow is denoted by arrows \textbackslash{}( I \textbackslash{}) and \textbackslash{}( i \textbackslash{}) in the circuit.
   - The direction of current is from the positive terminal of the battery through the ammeter and back to the negative terminal.

Overall, this illustration represents an experimental setup to study the photoelectric effect and measure the resulting current flow when light hits the photoelectric tube.

\#\# Dataset metadata

Quantum Phenomena; Physical Model Grounding Reasoning, Multi-Formula Reasoning

\paragraph{Recorded answer/context.}
A: A:  13.8kΩ

\paragraph{Figure/image path.}
/root/proposal\_for\_physic/science-mango/phyx\_mini\_run/phyx\_data/test\_image/612.png

\paragraph{Formalization target.}
create a compiling Lean file with sorry bodies at `PhyXMiniProblems/problem\_phyx\_mini\_0612.lean`. The Lean declarations must preserve the physical quantities, dimensions or dimensional roles, figure labels, governing-law hypotheses, and final relation expressed by this problem.
Use Mathlib/Physlib names found through LeanExplore where available. If a physics API is missing, introduce faithful local abstractions rather than scalar placeholder aliases.

\begin{theorem}[Physics formalization target]
\label{thm:physics:phyx_mini_0612:target}
\lean{PhyXMiniProblems.ProblemPhyXMini0612.problem_phyx_mini_0612}
\uses{def:physics:phyx-mini-0612:phyxminiproblems-problemphyxmini0612-photoelectriccircuitexperiment, def:physics:phyx-mini-0612:phyxminiproblems-problemphyxmini0612-matchesprimarycircuitfigure, def:physics:phyx-mini-0612:phyxminiproblems-problemphyxmini0612-matchesproblemscenario, def:physics:phyx-mini-0612:phyxminiproblems-problemphyxmini0612-usespreviouspartworkfunctionresult, def:physics:phyx-mini-0612:phyxminiproblems-problemphyxmini0612-usesstandardphysicalconstants, def:physics:phyx-mini-0612:phyxminiproblems-problemphyxmini0612-hasphysicalphotoelectriccircuitparameters, def:physics:phyx-mini-0612:phyxminiproblems-problemphyxmini0612-satisfiesphotoelectriccircuitlaws, def:physics:phyx-mini-0612:phyxminiproblems-problemphyxmini0612-isuniquematchingminimumstoppingresistance, def:physics:phyx-mini-0612:phyxminiproblems-problemphyxmini0612-recordeddatasetanswer}
The assigned autoformalize agent should translate this physics problem into Lean declarations in the covered file, with theorem and lemma proof bodies written as `by sorry`.
\end{theorem}
\begin{proof}
This is an autoformalization task, not a proof task. Produce statements that can later be proved by the physics prover without weakening the physical model.
\end{proof}

% --- Archon declaration topology begin ---
\section{Declaration topology}
\begin{definition}[electric Potential Dimension]
  \label{def:physics:phyx-mini-0612:phyxminiproblems-problemphyxmini0612-electricpotentialdimension}
  \lean{PhyXMiniProblems.ProblemPhyXMini0612.electricPotentialDimension}
  Electric potential has the dimension energy per charge
\end{definition}
\begin{proof}
  This is the declared model definition of electric Potential Dimension.
\end{proof}
\begin{definition}[electric Current Dimension]
  \label{def:physics:phyx-mini-0612:phyxminiproblems-problemphyxmini0612-electriccurrentdimension}
  \lean{PhyXMiniProblems.ProblemPhyXMini0612.electricCurrentDimension}
  Electric current has the dimension charge per time
\end{definition}
\begin{proof}
  This is the declared model definition of electric Current Dimension.
\end{proof}
\begin{definition}[electric Resistance Dimension]
  \label{def:physics:phyx-mini-0612:phyxminiproblems-problemphyxmini0612-electricresistancedimension}
  \lean{PhyXMiniProblems.ProblemPhyXMini0612.electricResistanceDimension}
  Electric resistance has the dimension potential divided by current
\end{definition}
\begin{proof}
  This is the declared model definition of electric Resistance Dimension.
\end{proof}
\begin{definition}[action Dimension]
  \label{def:physics:phyx-mini-0612:phyxminiproblems-problemphyxmini0612-actiondimension}
  \lean{PhyXMiniProblems.ProblemPhyXMini0612.actionDimension}
  Action, the dimensional role of Planck's constant, is energy times time
\end{definition}
\begin{proof}
  This is the declared model definition of action Dimension.
\end{proof}
\begin{definition}[Wavelength Quantity]
  \label{def:physics:phyx-mini-0612:phyxminiproblems-problemphyxmini0612-wavelengthquantity}
  \lean{PhyXMiniProblems.ProblemPhyXMini0612.WavelengthQuantity}
  A nonnegative physical wavelength
\end{definition}
\begin{proof}
  This is the declared model definition of Wavelength Quantity.
\end{proof}
\begin{definition}[Electric Potential Magnitude]
  \label{def:physics:phyx-mini-0612:phyxminiproblems-problemphyxmini0612-electricpotentialmagnitude}
  \lean{PhyXMiniProblems.ProblemPhyXMini0612.ElectricPotentialMagnitude}
  \uses{def:physics:phyx-mini-0612:phyxminiproblems-problemphyxmini0612-electricpotentialdimension}
  A nonnegative magnitude of electric potential
\end{definition}
\begin{proof}
  This is the declared model definition of Electric Potential Magnitude.
\end{proof}
\begin{definition}[Electric Current Magnitude]
  \label{def:physics:phyx-mini-0612:phyxminiproblems-problemphyxmini0612-electriccurrentmagnitude}
  \lean{PhyXMiniProblems.ProblemPhyXMini0612.ElectricCurrentMagnitude}
  \uses{def:physics:phyx-mini-0612:phyxminiproblems-problemphyxmini0612-electriccurrentdimension}
  A nonnegative magnitude of electric current
\end{definition}
\begin{proof}
  This is the declared model definition of Electric Current Magnitude.
\end{proof}
\begin{definition}[Electric Resistance Quantity]
  \label{def:physics:phyx-mini-0612:phyxminiproblems-problemphyxmini0612-electricresistancequantity}
  \lean{PhyXMiniProblems.ProblemPhyXMini0612.ElectricResistanceQuantity}
  \uses{def:physics:phyx-mini-0612:phyxminiproblems-problemphyxmini0612-electricresistancedimension}
  A nonnegative physical resistance
\end{definition}
\begin{proof}
  This is the declared model definition of Electric Resistance Quantity.
\end{proof}
\begin{definition}[Charge Magnitude Quantity]
  \label{def:physics:phyx-mini-0612:phyxminiproblems-problemphyxmini0612-chargemagnitudequantity}
  \lean{PhyXMiniProblems.ProblemPhyXMini0612.ChargeMagnitudeQuantity}
  A nonnegative magnitude of electric charge
\end{definition}
\begin{proof}
  This is the declared model definition of Charge Magnitude Quantity.
\end{proof}
\begin{definition}[Speed Quantity]
  \label{def:physics:phyx-mini-0612:phyxminiproblems-problemphyxmini0612-speedquantity}
  \lean{PhyXMiniProblems.ProblemPhyXMini0612.SpeedQuantity}
  A nonnegative physical speed
\end{definition}
\begin{proof}
  This is the declared model definition of Speed Quantity.
\end{proof}
\begin{definition}[Planck Constant Quantity]
  \label{def:physics:phyx-mini-0612:phyxminiproblems-problemphyxmini0612-planckconstantquantity}
  \lean{PhyXMiniProblems.ProblemPhyXMini0612.PlanckConstantQuantity}
  \uses{def:physics:phyx-mini-0612:phyxminiproblems-problemphyxmini0612-actiondimension}
  A nonnegative physical quantity with the dimension of action
\end{definition}
\begin{proof}
  This is the declared model definition of Planck Constant Quantity.
\end{proof}
\begin{definition}[nonnegative SIReadout]
  \label{def:physics:phyx-mini-0612:phyxminiproblems-problemphyxmini0612-nonnegativesireadout}
  \lean{PhyXMiniProblems.ProblemPhyXMini0612.nonnegativeSIReadout}
  Coherent-SI scalar component of a nonnegative dimensionful quantity
\end{definition}
\begin{proof}
  This is the declared model definition of nonnegative SIReadout.
\end{proof}
\begin{definition}[wavelength In Metres]
  \label{def:physics:phyx-mini-0612:phyxminiproblems-problemphyxmini0612-wavelengthinmetres}
  \lean{PhyXMiniProblems.ProblemPhyXMini0612.wavelengthInMetres}
  \uses{def:physics:phyx-mini-0612:phyxminiproblems-problemphyxmini0612-wavelengthquantity, def:physics:phyx-mini-0612:phyxminiproblems-problemphyxmini0612-nonnegativesireadout}
  Wavelength readout in metres
\end{definition}
\begin{proof}
  This is the declared model definition of wavelength In Metres.
\end{proof}
\begin{definition}[potential In Volts]
  \label{def:physics:phyx-mini-0612:phyxminiproblems-problemphyxmini0612-potentialinvolts}
  \lean{PhyXMiniProblems.ProblemPhyXMini0612.potentialInVolts}
  \uses{def:physics:phyx-mini-0612:phyxminiproblems-problemphyxmini0612-electricpotentialmagnitude, def:physics:phyx-mini-0612:phyxminiproblems-problemphyxmini0612-nonnegativesireadout}
  Electric-potential magnitude readout in volts
\end{definition}
\begin{proof}
  This is the declared model definition of potential In Volts.
\end{proof}
\begin{definition}[current In Amperes]
  \label{def:physics:phyx-mini-0612:phyxminiproblems-problemphyxmini0612-currentinamperes}
  \lean{PhyXMiniProblems.ProblemPhyXMini0612.currentInAmperes}
  \uses{def:physics:phyx-mini-0612:phyxminiproblems-problemphyxmini0612-electriccurrentmagnitude, def:physics:phyx-mini-0612:phyxminiproblems-problemphyxmini0612-nonnegativesireadout}
  Electric-current magnitude readout in amperes
\end{definition}
\begin{proof}
  This is the declared model definition of current In Amperes.
\end{proof}
\begin{definition}[resistance In Ohms]
  \label{def:physics:phyx-mini-0612:phyxminiproblems-problemphyxmini0612-resistanceinohms}
  \lean{PhyXMiniProblems.ProblemPhyXMini0612.resistanceInOhms}
  \uses{def:physics:phyx-mini-0612:phyxminiproblems-problemphyxmini0612-electricresistancequantity, def:physics:phyx-mini-0612:phyxminiproblems-problemphyxmini0612-nonnegativesireadout}
  Electric-resistance readout in ohms
\end{definition}
\begin{proof}
  This is the declared model definition of resistance In Ohms.
\end{proof}
\begin{definition}[resistance In Kilo Ohms]
  \label{def:physics:phyx-mini-0612:phyxminiproblems-problemphyxmini0612-resistanceinkiloohms}
  \lean{PhyXMiniProblems.ProblemPhyXMini0612.resistanceInKiloOhms}
  \uses{def:physics:phyx-mini-0612:phyxminiproblems-problemphyxmini0612-electricresistancequantity, def:physics:phyx-mini-0612:phyxminiproblems-problemphyxmini0612-resistanceinohms}
  Electric-resistance readout in kilohms
\end{definition}
\begin{proof}
  This is the declared model definition of resistance In Kilo Ohms.
\end{proof}
\begin{definition}[charge Magnitude In Coulombs]
  \label{def:physics:phyx-mini-0612:phyxminiproblems-problemphyxmini0612-chargemagnitudeincoulombs}
  \lean{PhyXMiniProblems.ProblemPhyXMini0612.chargeMagnitudeInCoulombs}
  \uses{def:physics:phyx-mini-0612:phyxminiproblems-problemphyxmini0612-chargemagnitudequantity, def:physics:phyx-mini-0612:phyxminiproblems-problemphyxmini0612-nonnegativesireadout}
  Electric-charge magnitude readout in coulombs
\end{definition}
\begin{proof}
  This is the declared model definition of charge Magnitude In Coulombs.
\end{proof}
\begin{definition}[speed In Metres Per Second]
  \label{def:physics:phyx-mini-0612:phyxminiproblems-problemphyxmini0612-speedinmetrespersecond}
  \lean{PhyXMiniProblems.ProblemPhyXMini0612.speedInMetresPerSecond}
  \uses{def:physics:phyx-mini-0612:phyxminiproblems-problemphyxmini0612-speedquantity, def:physics:phyx-mini-0612:phyxminiproblems-problemphyxmini0612-nonnegativesireadout}
  Speed readout in metres per second
\end{definition}
\begin{proof}
  This is the declared model definition of speed In Metres Per Second.
\end{proof}
\begin{definition}[planck Constant In Joule Seconds]
  \label{def:physics:phyx-mini-0612:phyxminiproblems-problemphyxmini0612-planckconstantinjouleseconds}
  \lean{PhyXMiniProblems.ProblemPhyXMini0612.planckConstantInJouleSeconds}
  \uses{def:physics:phyx-mini-0612:phyxminiproblems-problemphyxmini0612-planckconstantquantity, def:physics:phyx-mini-0612:phyxminiproblems-problemphyxmini0612-nonnegativesireadout}
  Planck-constant readout in joule-seconds
\end{definition}
\begin{proof}
  This is the declared model definition of planck Constant In Joule Seconds.
\end{proof}
\begin{definition}[energy In Joules]
  \label{def:physics:phyx-mini-0612:phyxminiproblems-problemphyxmini0612-energyinjoules}
  \lean{PhyXMiniProblems.ProblemPhyXMini0612.energyInJoules}
  Energy readout in joules
\end{definition}
\begin{proof}
  This is the declared model definition of energy In Joules.
\end{proof}
\begin{definition}[physlib Elementary Charge In Coulombs]
  \label{def:physics:phyx-mini-0612:phyxminiproblems-problemphyxmini0612-physlibelementarychargeincoulombs}
  \lean{PhyXMiniProblems.ProblemPhyXMini0612.physlibElementaryChargeInCoulombs}
  Physlib's exact elementary-charge magnitude, expressed in coulombs
\end{definition}
\begin{proof}
  This is the declared model definition of physlib Elementary Charge In Coulombs.
\end{proof}
\begin{definition}[physlib Speed Of Light In Metres Per Second]
  \label{def:physics:phyx-mini-0612:phyxminiproblems-problemphyxmini0612-physlibspeedoflightinmetrespersecond}
  \lean{PhyXMiniProblems.ProblemPhyXMini0612.physlibSpeedOfLightInMetresPerSecond}
  Physlib's speed of light, expressed in metres per second
\end{definition}
\begin{proof}
  This is the declared model definition of physlib Speed Of Light In Metres Per Second.
\end{proof}
\begin{definition}[physlib Planck Constant In Joule Seconds]
  \label{def:physics:phyx-mini-0612:phyxminiproblems-problemphyxmini0612-physlibplanckconstantinjouleseconds}
  \lean{PhyXMiniProblems.ProblemPhyXMini0612.physlibPlanckConstantInJouleSeconds}
  Physlib supplies the reduced Planck constant ℏ as a positive scalar in joule-seconds. The ordinary Planck constant used by the photon-energy law is h = 2πℏ
\end{definition}
\begin{proof}
  This is the declared model definition of physlib Planck Constant In Joule Seconds.
\end{proof}
\begin{definition}[nanometre In Metres]
  \label{def:physics:phyx-mini-0612:phyxminiproblems-problemphyxmini0612-nanometreinmetres}
  \lean{PhyXMiniProblems.ProblemPhyXMini0612.nanometreInMetres}
  One nanometre expressed in metres
\end{definition}
\begin{proof}
  This is the declared model definition of nanometre In Metres.
\end{proof}
\begin{definition}[Switch Position]
  \label{def:physics:phyx-mini-0612:phyxminiproblems-problemphyxmini0612-switchposition}
  \lean{PhyXMiniProblems.ProblemPhyXMini0612.SwitchPosition}
  The two labeled switch positions
\end{definition}
\begin{proof}
  This is the declared model definition of Switch Position.
\end{proof}
\begin{definition}[Figure Battery]
  \label{def:physics:phyx-mini-0612:phyxminiproblems-problemphyxmini0612-figurebattery}
  \lean{PhyXMiniProblems.ProblemPhyXMini0612.FigureBattery}
  The upper and lower 15.0 V cells in the figure
\end{definition}
\begin{proof}
  This is the declared model definition of Figure Battery.
\end{proof}
\begin{definition}[Circuit Side]
  \label{def:physics:phyx-mini-0612:phyxminiproblems-problemphyxmini0612-circuitside}
  \lean{PhyXMiniProblems.ProblemPhyXMini0612.CircuitSide}
  Left and right sides of the circuit drawing
\end{definition}
\begin{proof}
  This is the declared model definition of Circuit Side.
\end{proof}
\begin{definition}[Electrode]
  \label{def:physics:phyx-mini-0612:phyxminiproblems-problemphyxmini0612-electrode}
  \lean{PhyXMiniProblems.ProblemPhyXMini0612.Electrode}
  Cathode and anode of the photoelectric tube
\end{definition}
\begin{proof}
  This is the declared model definition of Electrode.
\end{proof}
\begin{definition}[Potential Sign]
  \label{def:physics:phyx-mini-0612:phyxminiproblems-problemphyxmini0612-potentialsign}
  \lean{PhyXMiniProblems.ProblemPhyXMini0612.PotentialSign}
  Signs printed under the electrodes in the V\_AC annotation
\end{definition}
\begin{proof}
  This is the declared model definition of Potential Sign.
\end{proof}
\begin{definition}[Figure Current]
  \label{def:physics:phyx-mini-0612:phyxminiproblems-problemphyxmini0612-figurecurrent}
  \lean{PhyXMiniProblems.ProblemPhyXMini0612.FigureCurrent}
  The two current arrows labeled in the diagram
\end{definition}
\begin{proof}
  This is the declared model definition of Figure Current.
\end{proof}
\begin{definition}[Figure Arrow Direction]
  \label{def:physics:phyx-mini-0612:phyxminiproblems-problemphyxmini0612-figurearrowdirection}
  \lean{PhyXMiniProblems.ProblemPhyXMini0612.FigureArrowDirection}
  Directions of the current arrows on the page
\end{definition}
\begin{proof}
  This is the declared model definition of Figure Arrow Direction.
\end{proof}
\begin{definition}[Photoelectric Circuit Figure]
  \label{def:physics:phyx-mini-0612:phyxminiproblems-problemphyxmini0612-photoelectriccircuitfigure}
  \lean{PhyXMiniProblems.ProblemPhyXMini0612.PhotoelectricCircuitFigure}
  \uses{def:physics:phyx-mini-0612:phyxminiproblems-problemphyxmini0612-electricpotentialmagnitude, def:physics:phyx-mini-0612:phyxminiproblems-problemphyxmini0612-electricresistancequantity, def:physics:phyx-mini-0612:phyxminiproblems-problemphyxmini0612-switchposition, def:physics:phyx-mini-0612:phyxminiproblems-problemphyxmini0612-figurebattery, def:physics:phyx-mini-0612:phyxminiproblems-problemphyxmini0612-circuitside, def:physics:phyx-mini-0612:phyxminiproblems-problemphyxmini0612-electrode, def:physics:phyx-mini-0612:phyxminiproblems-problemphyxmini0612-potentialsign, def:physics:phyx-mini-0612:phyxminiproblems-problemphyxmini0612-figurecurrent, def:physics:phyx-mini-0612:phyxminiproblems-problemphyxmini0612-figurearrowdirection}
  Typed presentation data from image 612. Physical component values remain dimensionful; the Booleans record literal schematic labels and symbols
\end{definition}
\begin{proof}
  This is the declared model definition of Photoelectric Circuit Figure.
\end{proof}
\begin{definition}[Photoelectric Circuit Experiment]
  \label{def:physics:phyx-mini-0612:phyxminiproblems-problemphyxmini0612-photoelectriccircuitexperiment}
  \lean{PhyXMiniProblems.ProblemPhyXMini0612.PhotoelectricCircuitExperiment}
  \uses{def:physics:phyx-mini-0612:phyxminiproblems-problemphyxmini0612-wavelengthquantity, def:physics:phyx-mini-0612:phyxminiproblems-problemphyxmini0612-electricpotentialmagnitude, def:physics:phyx-mini-0612:phyxminiproblems-problemphyxmini0612-electriccurrentmagnitude, def:physics:phyx-mini-0612:phyxminiproblems-problemphyxmini0612-electricresistancequantity, def:physics:phyx-mini-0612:phyxminiproblems-problemphyxmini0612-chargemagnitudequantity, def:physics:phyx-mini-0612:phyxminiproblems-problemphyxmini0612-speedquantity, def:physics:phyx-mini-0612:phyxminiproblems-problemphyxmini0612-planckconstantquantity, def:physics:phyx-mini-0612:phyxminiproblems-problemphyxmini0612-switchposition, def:physics:phyx-mini-0612:phyxminiproblems-problemphyxmini0612-photoelectriccircuitfigure}
  The experiment carries independent physical quantities and response maps. In particular, neither a stopping-resistance value nor an answer label is a field. photocurrentAt represents the ammeter magnitude for a switch position, wavelength, and candidate setting of the variable resistor R
\end{definition}
\begin{proof}
  This is the declared model definition of Photoelectric Circuit Experiment.
\end{proof}
\begin{definition}[Matches Primary Circuit Figure]
  \label{def:physics:phyx-mini-0612:phyxminiproblems-problemphyxmini0612-matchesprimarycircuitfigure}
  \lean{PhyXMiniProblems.ProblemPhyXMini0612.MatchesPrimaryCircuitFigure}
  \uses{def:physics:phyx-mini-0612:phyxminiproblems-problemphyxmini0612-potentialinvolts, def:physics:phyx-mini-0612:phyxminiproblems-problemphyxmini0612-resistanceinohms, def:physics:phyx-mini-0612:phyxminiproblems-problemphyxmini0612-photoelectriccircuitexperiment}
  Literal component values, polarities, labels, and arrow directions visible in the primary image. The opposite battery polarities here follow the raster, not the auxiliary caption's claim that the sources are simply in series
\end{definition}
\begin{proof}
  This is the declared model definition of Matches Primary Circuit Figure.
\end{proof}
\begin{definition}[Matches Problem Scenario]
  \label{def:physics:phyx-mini-0612:phyxminiproblems-problemphyxmini0612-matchesproblemscenario}
  \lean{PhyXMiniProblems.ProblemPhyXMini0612.MatchesProblemScenario}
  \uses{def:physics:phyx-mini-0612:phyxminiproblems-problemphyxmini0612-wavelengthinmetres, def:physics:phyx-mini-0612:phyxminiproblems-problemphyxmini0612-currentinamperes, def:physics:phyx-mini-0612:phyxminiproblems-problemphyxmini0612-nanometreinmetres, def:physics:phyx-mini-0612:phyxminiproblems-problemphyxmini0612-photoelectriccircuitexperiment}
  The prose readouts and observation at switch position a. The comparison of the two ammeter readouts records the stated current increase but prescribes no stopping resistance
\end{definition}
\begin{proof}
  This is the declared model definition of Matches Problem Scenario.
\end{proof}
\begin{definition}[Uses Previous Part Work Function Result]
  \label{def:physics:phyx-mini-0612:phyxminiproblems-problemphyxmini0612-usespreviouspartworkfunctionresult}
  \lean{PhyXMiniProblems.ProblemPhyXMini0612.UsesPreviousPartWorkFunctionResult}
  \uses{def:physics:phyx-mini-0612:phyxminiproblems-problemphyxmini0612-energyinjoules, def:physics:phyx-mini-0612:phyxminiproblems-problemphyxmini0612-physlibelementarychargeincoulombs, def:physics:phyx-mini-0612:phyxminiproblems-problemphyxmini0612-photoelectriccircuitexperiment}
  The prior part's material calibration, kept separate because its value is not printed in the supplied excerpt. This is a work-function result, not the current question's stopping-resistance result
\end{definition}
\begin{proof}
  This is the declared model definition of Uses Previous Part Work Function Result.
\end{proof}
\begin{definition}[Uses Standard Physical Constants]
  \label{def:physics:phyx-mini-0612:phyxminiproblems-problemphyxmini0612-usesstandardphysicalconstants}
  \lean{PhyXMiniProblems.ProblemPhyXMini0612.UsesStandardPhysicalConstants}
  \uses{def:physics:phyx-mini-0612:phyxminiproblems-problemphyxmini0612-chargemagnitudeincoulombs, def:physics:phyx-mini-0612:phyxminiproblems-problemphyxmini0612-speedinmetrespersecond, def:physics:phyx-mini-0612:phyxminiproblems-problemphyxmini0612-planckconstantinjouleseconds, def:physics:phyx-mini-0612:phyxminiproblems-problemphyxmini0612-physlibelementarychargeincoulombs, def:physics:phyx-mini-0612:phyxminiproblems-problemphyxmini0612-physlibspeedoflightinmetrespersecond, def:physics:phyx-mini-0612:phyxminiproblems-problemphyxmini0612-physlibplanckconstantinjouleseconds, def:physics:phyx-mini-0612:phyxminiproblems-problemphyxmini0612-photoelectriccircuitexperiment}
  The experiment uses the standard Physlib constants
\end{definition}
\begin{proof}
  This is the declared model definition of Uses Standard Physical Constants.
\end{proof}
\begin{definition}[Has Physical Photoelectric Circuit Parameters]
  \label{def:physics:phyx-mini-0612:phyxminiproblems-problemphyxmini0612-hasphysicalphotoelectriccircuitparameters}
  \lean{PhyXMiniProblems.ProblemPhyXMini0612.HasPhysicalPhotoelectricCircuitParameters}
  \uses{def:physics:phyx-mini-0612:phyxminiproblems-problemphyxmini0612-wavelengthinmetres, def:physics:phyx-mini-0612:phyxminiproblems-problemphyxmini0612-potentialinvolts, def:physics:phyx-mini-0612:phyxminiproblems-problemphyxmini0612-resistanceinohms, def:physics:phyx-mini-0612:phyxminiproblems-problemphyxmini0612-chargemagnitudeincoulombs, def:physics:phyx-mini-0612:phyxminiproblems-problemphyxmini0612-speedinmetrespersecond, def:physics:phyx-mini-0612:phyxminiproblems-problemphyxmini0612-planckconstantinjouleseconds, def:physics:phyx-mini-0612:phyxminiproblems-problemphyxmini0612-energyinjoules, def:physics:phyx-mini-0612:phyxminiproblems-problemphyxmini0612-photoelectriccircuitexperiment}
  Positivity and emission-side conditions for the physical setup
\end{definition}
\begin{proof}
  This is the declared model definition of Has Physical Photoelectric Circuit Parameters.
\end{proof}
\begin{definition}[Satisfies Photoelectric Circuit Laws]
  \label{def:physics:phyx-mini-0612:phyxminiproblems-problemphyxmini0612-satisfiesphotoelectriccircuitlaws}
  \lean{PhyXMiniProblems.ProblemPhyXMini0612.SatisfiesPhotoelectricCircuitLaws}
  \uses{def:physics:phyx-mini-0612:phyxminiproblems-problemphyxmini0612-wavelengthinmetres, def:physics:phyx-mini-0612:phyxminiproblems-problemphyxmini0612-potentialinvolts, def:physics:phyx-mini-0612:phyxminiproblems-problemphyxmini0612-currentinamperes, def:physics:phyx-mini-0612:phyxminiproblems-problemphyxmini0612-resistanceinohms, def:physics:phyx-mini-0612:phyxminiproblems-problemphyxmini0612-chargemagnitudeincoulombs, def:physics:phyx-mini-0612:phyxminiproblems-problemphyxmini0612-speedinmetrespersecond, def:physics:phyx-mini-0612:phyxminiproblems-problemphyxmini0612-planckconstantinjouleseconds, def:physics:phyx-mini-0612:phyxminiproblems-problemphyxmini0612-energyinjoules, def:physics:phyx-mini-0612:phyxminiproblems-problemphyxmini0612-photoelectriccircuitexperiment}
  Uniform physical laws used in the calculation: photon energy is Eγ = h c / λ; Einstein's photoelectric balance is Kmax = Eγ - φ; at switch a, the retarding voltage is the voltage-divider readout; photocurrent is zero exactly when e Vret ≥ Kmax. All laws quantify over arbitrary wavelengths or resistance settings and none mentions 13.8 kΩ or a displayed answer label
\end{definition}
\begin{proof}
  This is the declared model definition of Satisfies Photoelectric Circuit Laws.
\end{proof}
\begin{definition}[Photocurrent Ceases At]
  \label{def:physics:phyx-mini-0612:phyxminiproblems-problemphyxmini0612-photocurrentceasesat}
  \lean{PhyXMiniProblems.ProblemPhyXMini0612.PhotocurrentCeasesAt}
  \uses{def:physics:phyx-mini-0612:phyxminiproblems-problemphyxmini0612-electricresistancequantity, def:physics:phyx-mini-0612:phyxminiproblems-problemphyxmini0612-currentinamperes, def:physics:phyx-mini-0612:phyxminiproblems-problemphyxmini0612-photoelectriccircuitexperiment}
  The ammeter reads zero at the new wavelength and switch position a
\end{definition}
\begin{proof}
  This is the declared model definition of Photocurrent Ceases At.
\end{proof}
\begin{definition}[Is Minimum Stopping Resistance]
  \label{def:physics:phyx-mini-0612:phyxminiproblems-problemphyxmini0612-isminimumstoppingresistance}
  \lean{PhyXMiniProblems.ProblemPhyXMini0612.IsMinimumStoppingResistance}
  \uses{def:physics:phyx-mini-0612:phyxminiproblems-problemphyxmini0612-electricresistancequantity, def:physics:phyx-mini-0612:phyxminiproblems-problemphyxmini0612-resistanceinohms, def:physics:phyx-mini-0612:phyxminiproblems-problemphyxmini0612-photoelectriccircuitexperiment, def:physics:phyx-mini-0612:phyxminiproblems-problemphyxmini0612-photocurrentceasesat}
  A physical resistance is a minimum stopping setting when it stops the current and no other stopping setting has a smaller coherent-SI resistance readout
\end{definition}
\begin{proof}
  This is the declared model definition of Is Minimum Stopping Resistance.
\end{proof}
\begin{definition}[Answer Choice]
  \label{def:physics:phyx-mini-0612:phyxminiproblems-problemphyxmini0612-answerchoice}
  \lean{PhyXMiniProblems.ProblemPhyXMini0612.AnswerChoice}
  Labels printed beside the four multiple-choice answers
\end{definition}
\begin{proof}
  This is the declared model definition of Answer Choice.
\end{proof}
\begin{definition}[displayed Resistance In Kilo Ohms]
  \label{def:physics:phyx-mini-0612:phyxminiproblems-problemphyxmini0612-displayedresistanceinkiloohms}
  \lean{PhyXMiniProblems.ProblemPhyXMini0612.displayedResistanceInKiloOhms}
  \uses{def:physics:phyx-mini-0612:phyxminiproblems-problemphyxmini0612-answerchoice}
  Resistance readout printed for each choice, in kilohms
\end{definition}
\begin{proof}
  This is the declared model definition of displayed Resistance In Kilo Ohms.
\end{proof}
\begin{definition}[displayed Resistance Tolerance In Kilo Ohms]
  \label{def:physics:phyx-mini-0612:phyxminiproblems-problemphyxmini0612-displayedresistancetoleranceinkiloohms}
  \lean{PhyXMiniProblems.ProblemPhyXMini0612.displayedResistanceToleranceInKiloOhms}
  Half of the 0.1 kΩ precision to which the choices are displayed
\end{definition}
\begin{proof}
  This is the declared model definition of displayed Resistance Tolerance In Kilo Ohms.
\end{proof}
\begin{definition}[Matches Displayed Minimum Stopping Resistance]
  \label{def:physics:phyx-mini-0612:phyxminiproblems-problemphyxmini0612-matchesdisplayedminimumstoppingresistance}
  \lean{PhyXMiniProblems.ProblemPhyXMini0612.MatchesDisplayedMinimumStoppingResistance}
  \uses{def:physics:phyx-mini-0612:phyxminiproblems-problemphyxmini0612-resistanceinkiloohms, def:physics:phyx-mini-0612:phyxminiproblems-problemphyxmini0612-photoelectriccircuitexperiment, def:physics:phyx-mini-0612:phyxminiproblems-problemphyxmini0612-isminimumstoppingresistance, def:physics:phyx-mini-0612:phyxminiproblems-problemphyxmini0612-answerchoice, def:physics:phyx-mini-0612:phyxminiproblems-problemphyxmini0612-displayedresistanceinkiloohms, def:physics:phyx-mini-0612:phyxminiproblems-problemphyxmini0612-displayedresistancetoleranceinkiloohms}
  A displayed choice agrees, to displayed precision, with a physical minimum
\end{definition}
\begin{proof}
  This is the declared model definition of Matches Displayed Minimum Stopping Resistance.
\end{proof}
\begin{definition}[Is Unique Matching Minimum Stopping Resistance]
  \label{def:physics:phyx-mini-0612:phyxminiproblems-problemphyxmini0612-isuniquematchingminimumstoppingresistance}
  \lean{PhyXMiniProblems.ProblemPhyXMini0612.IsUniqueMatchingMinimumStoppingResistance}
  \uses{def:physics:phyx-mini-0612:phyxminiproblems-problemphyxmini0612-photoelectriccircuitexperiment, def:physics:phyx-mini-0612:phyxminiproblems-problemphyxmini0612-answerchoice, def:physics:phyx-mini-0612:phyxminiproblems-problemphyxmini0612-matchesdisplayedminimumstoppingresistance}
  Exactly one displayed choice agrees with the physical minimum
\end{definition}
\begin{proof}
  This is the declared model definition of Is Unique Matching Minimum Stopping Resistance.
\end{proof}
\begin{definition}[recorded Dataset Answer]
  \label{def:physics:phyx-mini-0612:phyxminiproblems-problemphyxmini0612-recordeddatasetanswer}
  \lean{PhyXMiniProblems.ProblemPhyXMini0612.recordedDatasetAnswer}
  \uses{def:physics:phyx-mini-0612:phyxminiproblems-problemphyxmini0612-answerchoice}
  The answer label recorded in the source dataset
\end{definition}
\begin{proof}
  This is the declared model definition of recorded Dataset Answer.
\end{proof}
\begin{lemma}[minimum stopping Resistance satisfies threshold Equality]
  \label{lem:physics:phyx-mini-0612:phyxminiproblems-problemphyxmini0612-minimum-stoppingresistance-satisfies-thresholdequality}
  \lean{PhyXMiniProblems.ProblemPhyXMini0612.minimum_stoppingResistance_satisfies_thresholdEquality}
  \uses{def:physics:phyx-mini-0612:phyxminiproblems-problemphyxmini0612-electricresistancequantity, def:physics:phyx-mini-0612:phyxminiproblems-problemphyxmini0612-potentialinvolts, def:physics:phyx-mini-0612:phyxminiproblems-problemphyxmini0612-chargemagnitudeincoulombs, def:physics:phyx-mini-0612:phyxminiproblems-problemphyxmini0612-energyinjoules, def:physics:phyx-mini-0612:phyxminiproblems-problemphyxmini0612-photoelectriccircuitexperiment, def:physics:phyx-mini-0612:phyxminiproblems-problemphyxmini0612-hasphysicalphotoelectriccircuitparameters, def:physics:phyx-mini-0612:phyxminiproblems-problemphyxmini0612-satisfiesphotoelectriccircuitlaws, def:physics:phyx-mini-0612:phyxminiproblems-problemphyxmini0612-isminimumstoppingresistance}
  At a least stopping resistance, the voltage divider supplies exactly the stopping energy. This is a derived threshold relation, not an assumption
\end{lemma}
\begin{proof}
  The conclusion follows from the governing relations and figure readouts named in the statement of minimum stopping Resistance satisfies threshold Equality.
\end{proof}
\begin{lemma}[exists minimum stopping Resistance near thirteen Point Eight Kilo Ohms]
  \label{lem:physics:phyx-mini-0612:phyxminiproblems-problemphyxmini0612-exists-minimum-stoppingresistance-near-thirteenpointeightkiloohms}
  \lean{PhyXMiniProblems.ProblemPhyXMini0612.exists_minimum_stoppingResistance_near_thirteenPointEightKiloOhms}
  \uses{def:physics:phyx-mini-0612:phyxminiproblems-problemphyxmini0612-resistanceinkiloohms, def:physics:phyx-mini-0612:phyxminiproblems-problemphyxmini0612-photoelectriccircuitexperiment, def:physics:phyx-mini-0612:phyxminiproblems-problemphyxmini0612-matchesprimarycircuitfigure, def:physics:phyx-mini-0612:phyxminiproblems-problemphyxmini0612-matchesproblemscenario, def:physics:phyx-mini-0612:phyxminiproblems-problemphyxmini0612-usespreviouspartworkfunctionresult, def:physics:phyx-mini-0612:phyxminiproblems-problemphyxmini0612-usesstandardphysicalconstants, def:physics:phyx-mini-0612:phyxminiproblems-problemphyxmini0612-hasphysicalphotoelectriccircuitparameters, def:physics:phyx-mini-0612:phyxminiproblems-problemphyxmini0612-satisfiesphotoelectriccircuitlaws, def:physics:phyx-mini-0612:phyxminiproblems-problemphyxmini0612-isminimumstoppingresistance}
  The physical minimum exists and rounds to 13.8 kΩ after applying the previous work-function result, the 65.0 nm readout, the photoelectric laws, and the switch-a voltage divider
\end{lemma}
\begin{proof}
  The conclusion follows from the governing relations and figure readouts named in the statement of exists minimum stopping Resistance near thirteen Point Eight Kilo Ohms.
\end{proof}
% --- Archon declaration topology end ---
% --- Archon physics formalization source end ---
